%%%%%%%% ICML 2026 CAMERA-READY (non-anonymous) %%%%%%%%%%%%%%%%%
%% gDTR: Layer-wise Settling Depth Reveals Biological Grammar
%% in Genomic Foundation Models
%% Accepted at the ICML 2026 Workshop on Generative and Agentic AI
%% for Biology (GenBio) (short paper: 4 pages + appendix)

\documentclass{article}

% Recommended, but optional, packages for figures and better typesetting:
\usepackage{microtype}
\usepackage{afterpage}     % \afterpage{\clearpage} to flush a page early
\usepackage{graphicx}

% Relax LaTeX float-placement constraints so Fig.~\ref{fig:shallowness}
% lands at the top of the page where Section~3 begins (rather than
% being deferred a page later because of the default conservative
% topfraction).
\renewcommand{\topfraction}{0.92}
\renewcommand{\textfraction}{0.08}
\renewcommand{\floatpagefraction}{0.80}
\renewcommand{\dbltopfraction}{0.92}
\renewcommand{\dblfloatpagefraction}{0.80}
\usepackage{subcaption}
\usepackage{booktabs}      % professional tables
\usepackage{multirow}
\usepackage{array}
\usepackage{siunitx}       % aligned numerics in tables
\usepackage{xcolor}
\usepackage{tikz}
\usetikzlibrary{arrows.meta,shapes.misc}
% Make \sffamily resolve to DejaVu Sans, the EXACT face used by the
% matplotlib raster figures (Fig 2--6). Combined with the TikZ panels
% of Fig 1 also using \sffamily, every figure in the paper -- TikZ or
% matplotlib -- now renders sans-serif text in the same family and
% weight class.
\usepackage[scaled=0.95]{DejaVuSans}

% Keep float-only pages from vertically centering figures with large blank tops.
\makeatletter
\setlength{\@fptop}{0pt}
\setlength{\@fpsep}{10pt plus 1fil}
\setlength{\@dblfptop}{0pt}
\setlength{\@dblfpsep}{10pt plus 1fil}
\makeatother

% Float separations relaxed toward the LaTeX defaults (textfloatsep 20pt,
% floatsep/intextsep 12pt) for camera-ready compliance with the ICML rule
% against compressing vertical space. Kept mildly tightened (not fully at
% default) only to avoid a visible gap in the text column below Fig 3.
\setlength{\textfloatsep}{16pt plus 2pt minus 4pt}
\setlength{\floatsep}{11pt plus 2pt minus 2pt}
\setlength{\intextsep}{11pt plus 2pt minus 2pt}
\setlength{\dbltextfloatsep}{16pt plus 2pt minus 4pt}
\setlength{\dblfloatsep}{11pt plus 2pt minus 2pt}

\usepackage{hyperref}
\newcommand{\theHalgorithm}{\arabic{algorithm}}

% Use the following line for the initial blind version submitted for review:
% \usepackage{icml2026}

% For preprint, use
% \usepackage[preprint]{icml2026}

% If accepted, instead use the following line for the camera-ready submission:
\usepackage[accepted]{icml2026}

% Camera-ready footer: this paper appears at the GenBio workshop, NOT in the
% main ICML proceedings. Override the default ``appearing in'' note here so the
% shared icml2026.sty stays untouched. \ICML@appearing is expanded lazily at use
% time, so this preamble redefinition is picked up by the first-page footer.
\makeatletter
\renewcommand{\ICML@appearing}{\textit{Accepted at the 2026 Workshop on Generative and Agentic AI for Biology (ICML 2026)}}
\makeatother

\usepackage{amsmath}
\usepackage{amssymb}
\usepackage{mathtools}
\usepackage{amsthm}

% if you use cleveref..
\usepackage[capitalize,noabbrev]{cleveref}
\usepackage{placeins}   % \FloatBarrier for pinning figures within a section

% \paragraph spacing relaxed toward the LaTeX default before-skip (default is
% 3.25ex). Kept moderately tightened (1.8ex) so run-in headings (Evo 2's idle
% final block, Tuned-lens sanity check) stay readable without the aggressive
% compression that the ICML camera-ready guidelines discourage.
% Placed AFTER all package loads so it overrides any \paragraph from icml2026.sty.
\makeatletter
\renewcommand\paragraph{\@startsection{paragraph}{4}{\z@}%
  {1.8ex \@plus 0.4ex \@minus 0.2ex}%
  {-0.6em}%
{\normalfont\normalsize\bfseries}}
\makeatother

%%%%%%%%%%%%%%%%%%%%%%%%%%%%%%%%
% THEOREMS / DEFINITIONS
%%%%%%%%%%%%%%%%%%%%%%%%%%%%%%%%
\theoremstyle{plain}
\newtheorem{theorem}{Theorem}[section]
\newtheorem{proposition}[theorem]{Proposition}
\newtheorem{lemma}[theorem]{Lemma}
\newtheorem{corollary}[theorem]{Corollary}
\theoremstyle{definition}
\newtheorem{definition}[theorem]{Definition}
\newtheorem{assumption}[theorem]{Assumption}
\theoremstyle{remark}
\newtheorem{remark}[theorem]{Remark}

%%%%%%%%%%%%%%%%%%%%%%%%%%%%%%%%
% MACROS
%%%%%%%%%%%%%%%%%%%%%%%%%%%%%%%%
\newcommand{\gDTR}{\textsc{GDTR}}
\newcommand{\DTR}{\textsc{DTR}}
\newcommand{\hl}{h_{\ell}}
\newcommand{\hnorm}{h_{\mathrm{norm}}}
\newcommand{\Dcos}{D_{\cos}}
\newcommand{\Djsd}{D_{\mathrm{jsd}}}
\newcommand{\DDcos}{\Delta D_{\cos}}
\newcommand{\DDjsd}{\Delta D_{\mathrm{jsd}}}
\newcommand{\runmin}{\operatorname{run\text{-}min}}

% Running title (kept short for the page header)
\icmltitlerunning{GDTR: Layer-wise Settling Depth Reveals Biological Grammar}

\begin{document}

\twocolumn[
  \icmltitle{\gDTR{}: Layer-wise Settling Depth Reveals Biological Grammar in Genomic Foundation Models}

  % Camera-ready, non-anonymous: the [accepted] option is ON, so the style
  % file shows the real author names, affiliations, and corresponding author.
  % \icmlsetsymbol{equal}{*}

  \begin{icmlauthorlist}
    \icmlauthor{Yoonjin Cho}{med}
    \icmlauthor{Jiheon Kang}{ee}
    \icmlauthor{Subin Park}{cs}
    \icmlauthor{Sangwoo Kim}{bmsi}
  \end{icmlauthorlist}

  \icmlaffiliation{med}{Yonsei University College of Medicine, Seoul, Korea}
  \icmlaffiliation{bmsi}{Department of Biomedical Systems Informatics, Yonsei University College of Medicine, Seoul, Korea}
  \icmlaffiliation{ee}{Department of Electrical and Electronic Engineering, Yonsei University, Seoul, Korea}
  \icmlaffiliation{cs}{Department of Computer Science, Yonsei University, Seoul, Korea}
  \icmlcorrespondingauthor{Sangwoo Kim}{swkim@yuhs.ac}

  \icmlkeywords{Interpretability, Genomic Foundation Models,
    Logit Lens, Tuned Lens, Mutational Signatures,
  Variant Effect Prediction, Mechanistic Interpretability}

  \vskip 0.3in
]

% No special notice (required even if empty).
\printAffiliationsAndNotice{}

%%%%%%%%%%%%%%%%%%%%%%%%%%%%%%%%
% ABSTRACT
%%%%%%%%%%%%%%%%%%%%%%%%%%%%%%%%
\begin{abstract}
  Genomic foundation models capture sequence regularities, yet existing interpretability tools rarely ask \emph{where} in the layer stack a biological grammar becomes stable. We introduce \gDTR{} (\emph{Genomic Deep-Thinking Ratio}), a training-free residual-stream lens that assigns each nucleotide token a \emph{settling depth} $c(t)$: the first layer at which its representation stabilises against the post-final-norm reference. On Evo~2~7B, splice donor and acceptor sites settle approximately two layers earlier than intronic contexts, enhancer-like cCREs show a smaller but measurable shift, and a chr22 calibration transfers to held-out chr17. Perturbing canonical splice donors shows that the signal is bidirectional: disrupting the central GT motif deepens settling, whereas shuffling the flanking grammar makes the preserved motif settle earlier. Differential \gDTR{} further reveals consequence-associated peak-disruption depths across ClinVar variants, with synonymous substitutions peaking deepest but with broad class overlap. \gDTR{} therefore provides a layer-wise interpretability axis for genomic foundation models, complementary to existing prediction and variant-scoring tools.
\end{abstract}

%%%%%%%%%%%%%%%%%%%%%%%%%%%%%%%%
% Fig 1 — conceptual (top of paper)
%%%%%%%%%%%%%%%%%%%%%%%%%%%%%%%%
\begin{figure*}[!t]
  \centering
  \includegraphics[width=0.98\textwidth]{figures/fig_1_final.png}
  \caption{\gDTR{} measures the layer at which a token's representation settles into the model's output-ready frame. Each Evo~2 layer emits a residual-stream tap $h_\ell(t)$; the \emph{cosine lens} compares it to the post-final-norm reference $\hnorm$ via $\Dcos(\ell,t)\!=\!1\!-\!\cos(h_\ell(t),\hnorm(t))$ (\emph{distance to final state}). The \emph{running-minimum envelope} $m_\ell\!=\!\min_{j\le\ell}\Dcos(j)$ enforces a monotone non-increasing trace, and the settling depth $c(t)$ is the first (shallowest) layer at which $m_\ell\!\le\!\gamma$ (Eq.~\ref{eq:settling}). In the illustrative trace shown, $c(t)\!=\!29$: a raw distance at a later layer rises back above $\gamma$, but the envelope carries the smaller value at layer~29 forward, so every layer at or beyond layer~29 satisfies $m_\ell\!\le\!\gamma$ while shallower layers ($\ell\!\le\!28$) do not. Lower $c(t)$ means earlier stabilisation. The layer numbers in this schematic are pedagogical; Evo~2's actual layer-30 saturation is discussed in \S\ref{sec:framework} and App.~\ref{app:architectural-quirk}.}
  \label{fig:schematic}
\end{figure*}

%%%%%%%%%%%%%%%%%%%%%%%%%%%%%%%%
% 1. INTRODUCTION
%%%%%%%%%%%%%%%%%%%%%%%%%%%%%%%%
\section{Introduction}
\label{sec:intro}

Genomic foundation models~\cite{nguyen2026evo2,nguyen2023hyenadna,dallatorre2024nt,zhou2024dnabert2,schiff2024caduceus} are increasingly used across functional-genomics tasks, but their internal layer-wise dynamics remain difficult to interpret. Existing tools ask \emph{where} models attend~\cite{abnar2020rollout,avsec2021enformer}, \emph{what} features they encode~\cite{bricken2023monosemanticity,cunningham2023sparse}, or \emph{how} variants shift final representations~\cite{cheng2023alphamissense,jaganathan2019spliceai}. Even recent sparse-autoencoder analyses of Evo~2~\cite{nguyen2026evo2,deng2025interpreting} recover interpretable genomic features from selected layers, answering \emph{what} the model represents. A complementary question remains largely unmeasured: \emph{when, along the layer stack, does a biological sequence element become stable?} On Evo~2, a splice donor's representation stabilises about two layers earlier than a nearby intronic position, even though both end at the same final layer.

We introduce \gDTR{} (\emph{Genomic Deep-Thinking Ratio}), a training-free residual-stream lens (a per-layer readout of the model's hidden state) for this question. \gDTR{} assigns each nucleotide token a \emph{settling depth} $c(t)$: the first layer at which its representation stabilises against the model's output-ready reference. We use \emph{biological grammar} to mean the context-dependent integration of a local motif with its flanking sequence, as in splice donors where the central GT dinucleotide is functional only in context. \gDTR{} adapts the Deep-Thinking Ratio idea~\cite{chen2026deepthinking} to genomic causal LMs with a cosine residual-stream lens, a running-minimum envelope that smooths non-monotone layer-by-layer trajectories, and an Evo~2-specific post-final-norm reference. Fig.~\ref{fig:schematic} summarises the readout; details are in \S\ref{sec:framework}.

Our central claim is that settling depth is biologically structured. On Evo~2~7B, a single chr22 calibration of $c(t)$ transfers to held-out chr17; splice sites and enhancer-like cCREs settle earlier than surrounding contexts; and matched motif and flank perturbations on the same donor loci move settling depth in opposite directions, separating motif detection from flanking-context integration. Subtracting the reference-allele cosine trajectory from the alternate-allele one, the same lens shows consequence-associated peak-disruption depths across ClinVar variant classes. \gDTR{} therefore provides a layer-resolved interpretability axis for genomic foundation models, complementary to final prediction scores and variant scorers.

%%%%%%%%%%%%%%%%%%%%%%%%%%%%%%%%
% 2. FRAMEWORK
%%%%%%%%%%%%%%%%%%%%%%%%%%%%%%%%
\section{The \gDTR{} Framework}
\label{sec:framework}

For a nucleotide sequence of length $T$ processed by a causal language model with $L$ layers, we extract the residual-stream activation $\hl(t)$ at every layer $\ell\!\in\!\{1,\ldots,L\}$ and token position $t$. We replace the JSD lens of the parent DTR~\cite{chen2026deepthinking} with a cosine residual-stream lens that operates directly in hidden-state geometry,
\begin{equation}
  \Dcos(\ell, t) \;=\; 1 - \cos\!\bigl(\hl(t),\, \hnorm(t)\bigr),
  \label{eq:cos-lens}
\end{equation}
where $\hnorm(t)$ is the post-final-RMSNorm state that the model's unembedding head reads (layers indexed $0,\ldots,L\!-\!1$ throughout the Evo~2 analysis).
% REVIEW FIX (blue note at Eq.1 "or cos, we can change"): justify cosine
% over JSD via Evo-2's 1-mer vocabulary; last sentence reconciles with the
% JSD AUROC in App. C so the two sections do not contradict.
Evo~2 tokenises at single-nucleotide (1-mer) resolution, so its next-token mass concentrates on an approximately 4-symbol alphabet, and the parent DTR's logit-lens JSD collapses to a near-binary trace with little discriminative resolution. The cosine distance in the 4096-dimensional residual stream descends smoothly enough to threshold into $c(t)$ and avoids the unembedding projection, which is unreliable through Evo~2's rotations and pre-final saturation (App.~\ref{app:architectural-quirk}). We retain JSD only as a cross-lens check for variant discrimination (App.~\ref{app:auroc-detail}), where aggregation over positions restores its signal. We define the \emph{settling depth} via a running-minimum envelope,
\begin{equation}
  c(t) \;=\; \min\bigl\{\, \ell \;:\; \runmin\,\Dcos(\ell, t) \le \gamma_{\cos}\bigr\},
  \label{eq:settling}
\end{equation}
where $\runmin\,\Dcos(\ell,t) = \min_{k\le\ell}\Dcos(k,t)$. Note that this is equivalent to the first-passage time $\min\{\ell : D_{\cos}(\ell, t) \leq \gamma_{\cos}\}$; the running-minimum is written explicitly to reflect the monotone envelope used for visualization in Fig.~\ref{fig:schematic}.

Lower $c(t)$ means earlier stabilisation; higher $c(t)$ means that the token remains unresolved until later layers. NLP transformers exhibit near-monotonic logit-lens convergence~\cite{nostalgebraist2020logitlens,belrose2023tunedlens,pal2023futurelens}; genomic CLMs do not (Hyena alignment rotations, Evo~2 pre-final saturation), so the running-min reduces this non-monotonicity to a monotone-non-increasing trace. The threshold $\gamma_{\cos}$ is set by \emph{regional $q_{70}$ calibration}: the 70th percentile of the running-min at the penultimate layer over the analysis region. The operating point sits within a wide, flat plateau identified by a $5\!\times\!5$ grid search over $(\gamma_{\cos}, \rho)$ (App.~\ref{app:gamma-sweep}).

\paragraph{Settling depth is two-sided.} A lower $c(t)$ does not by itself imply a stronger biological signal. A token can settle early because (a) its representation is constrained by a biological grammar that the model commits to early, \emph{or} (b) its surrounding context is simpler, so the running-min envelope reaches $\gamma_{\cos}$ with little integration. We exploit this two-sidedness directly: the perturbation experiments in \S\ref{sec:motif-controls} are designed to push $c(t)$ in opposite directions, separately stressing motif detection and flanking-grammar integration. Throughout the paper we therefore interpret depth signatures as \emph{bidirectional}. Both directions of movement can carry information about motif detection and context integration, but they should not be read as a monotone scale of biological signal strength.
% REVIEW FIX (concern #1: Table A11 looks like an internal contradiction):
% pre-frame the non-canonical-donor ordering as expected under arm (b) so
% the headline bidirectional claim covers it; full reconciliation in App. D.2.
As a concrete instance of arm (b), non-canonical splice donors settle \emph{earlier} than the dominant GT-AG class (App.~\ref{app:splice-canonical}), reflecting tighter flanking constraint rather than a stronger splice signal.

\paragraph{Representational saturation before the final layer.} The reference $\hnorm$ is fixed by the model as the post-final-RMSNorm state that the unembedding reads, independently of any tap we select. Measured against it, Evo~2 reaches the post-norm direction before the last attention layer executes: on chr22 control sequences, $\max_{t}|h_{30}(t)\!-\!h_{31}(t)|\!=\!0$ exactly (layer~31 is a residual passthrough), while $\cos(h_{29},\hnorm)\!=\!-0.013$ versus $\cos(h_{30},\hnorm)\!=\!\cos(h_{31},\hnorm)\!=\!+0.6855$. Layer~30 is thus the step that rotates the representation into the output-ready frame, leaving layer~29 as the deepest tap not yet aligned with it. Since the running-min envelope (Eq.~\ref{eq:settling}) records how early a token reaches $\hnorm$, the discriminating variation in settling depth lies in these pre-rotation layers. We report $L^\star\!=\!29$ as that empirical boundary rather than as a claim about the model's architectural design intent; full evidence is in App.~\ref{app:architectural-quirk}.

\paragraph{Tuned-lens consistency check against frame mismatch.} Frame mismatch is the worry that interior layers sit in a different linear basis from $\hnorm$, so a raw cosine distance would track basis rotation rather than genuine convergence. A per-layer tuned-lens affine fit~\cite{belrose2023tunedlens} reaches $\ge 98\%$ MSE recovery on $30/32$ Evo~2 layers (canonical $L^\star\!=\!29$: $0.9996$; App.~\ref{app:tuned-lens}), so the interior taps are one affine transform away from the output frame. Gross frame mismatch is therefore unlikely to drive the \gDTR{} signal. No affine weights are used at inference, so the framework remains training-free.

% Force Fig.~\ref{fig:shallowness} to land at the TOP of the page
% where Section~3 begins. The figure* source is placed BEFORE the
% \section so the [!t] float can occupy the very top of that page,
% with \afterpage{\clearpage} flushing the previous page early.
%\afterpage{\clearpage}

\begin{figure*}[!t]
  \centering
  \includegraphics[width=0.95\textwidth]{figures/fig_2_shallowness.png}
  \caption{Splice and enhancer annotations shift toward shallower settling. Both panels report Cohen's $d$ against the pooled chr17+chr22 intron baseline ($\overline{c}\!=\!27.72$, red dashed). \textbf{(a)} Per-context mean settling depth $\overline{c}$ on canonical genomic contexts (frozen $\gamma_{\cos}\!=\!0.397$). Splice donor and acceptor (blue) sit approximately 2 layers below the intron baseline. The asterisk on ``5$'$~UTR$^*$'' marks an entropy-coupled context (\S\ref{sec:discussion}) reported separately from the entropy-controlled hierarchy. Right of each bar: Cohen's $d$ versus the intron baseline (Mann--Whitney significance: $\,*$, $**$, $***$, $****$ at $0.05$, $0.01$, $10^{-3}$, $10^{-4}$). \textbf{(b)} Cohen's $d$ for regulatory annotations: splice donor (reference; $-0.354$), ENCODE cCRE-ELS ($-0.118$), GTEx Whole-Blood eQTL ($+0.064$), GWAS Catalog ($+0.021$). cCRE-ELS yields a measurable enhancer signal; SNP-level eQTL and GWAS shifts are negligible on the effect-size scale.}
  \label{fig:shallowness}
\end{figure*}

%%%%%%%%%%%%%%%%%%%%%%%%%%%%%%%%
% 3. BIOLOGICAL UTILITY
%%%%%%%%%%%%%%%%%%%%%%%%%%%%%%%%
\section{Results}
\label{sec:utility}

\paragraph{Statistical conventions.} The test in each subsection is fixed by the structure of its data rather than chosen for convenience: \S\ref{sec:shallowness} uses the Mann--Whitney $U$ test (independent two-sample comparisons at large $n$), \S\ref{sec:motif-controls} uses the paired Wilcoxon signed-rank test (the same donor loci under matched perturbations), and \S\ref{sec:variant-disruption} uses a Kruskal--Wallis omnibus with Bonferroni-corrected Dunn post-hoc tests (three or more unmatched groups of discrete, tied argmax layers). We report Cohen's $d$ wherever comparisons are pairwise, $\varepsilon^2$ for the Kruskal--Wallis omnibus, and the rank-biserial $r$ for Dunn contrasts, so that effect magnitudes remain comparable across subsections.

\subsection{Splice and Enhancer Annotations Settle Early}
\label{sec:shallowness}

We first test whether $c(t)$ produces a transferable readout at the genome scale. Chromosome~22 is the calibration set ($\gamma_{\cos}\!=\!0.397$ from $q_{70}$ at the penultimate layer); chromosome~17 is held out for validation with that $\gamma$ frozen. The chr17 $q_{70}$ recomputed independently equals $0.394$ ($|\Delta|\!=\!0.0023$); the chr17 splice-donor effect reaches $94.6\%$ of the chr22 magnitude (Cohen's $d\!=\!-0.349$ vs intron), and the relative ordering of the seven canonical contexts is preserved (Spearman $\rho\!=\!0.93$). The calibration transfers without re-tuning, and we use this single $\gamma$ for all subsequent context- and variant-level comparisons. All context-level summaries below report the per-position $c(t)$ pooled over the analysis windows, summarised by Cohen's $d$ versus intron, which normalises by the pooled within-context standard deviation and so absorbs within-region variability by construction.

\textbf{Context-level ordering.} On canonical genomic contexts (Fig.~\ref{fig:shallowness}a), splice donor and acceptor sites settle approximately 2 layers earlier than the intronic baseline, and the remaining contexts cluster near or above it. Given the large number of pooled positions ($\approx 10^7$), we treat effect size rather than nominal significance as the primary evidence and Mann--Whitney $p$-values as a secondary check. Among the regulatory annotations (Fig.~\ref{fig:shallowness}b), all Cohen's $d$ values are computed against the same pooled chr17+chr22 intron baseline ($\overline{c}\!=\!27.72$) used in panel~(a), so cross-panel magnitudes are directly comparable. ENCODE cCRE-ELS enhancer-like elements~\cite{moore2020encode} are the clearest non-splice signal at Cohen's $d\!=\!-0.118$ on this baseline (vs.\ splice-donor $d\!=\!-0.354$): biologically modest but statistically robust, and the effect transfers from chr22 to chr17. SNP-level annotations (GTEx Whole-Blood chr22 eQTLs~\cite{gtex2020}: $n\!=\!32{,}456$, $d\!=\!+0.064$; GWAS Catalog chr22 SNPs~\cite{sollis2023gwas}: $n\!=\!6{,}900$, $d\!=\!+0.021$) yield negligible effect sizes that we do not interpret further. The 5$'$~UTR shift is reported separately because it showed anomalous entropy coupling on the 120-window control panel ($\rho\!=\!+0.41$), although this coupling attenuates substantially at full-chr22 scale ($\rho\!=\!+0.054$, with $|\rho|\!\le\!0.224$ across all seven canonical contexts; App.~\ref{app:restriction-of-range}). We keep 5$'$~UTR on a separate axis, a conservative carryover from the smaller panel.

\textbf{Entropy does not explain the splice signal.} A natural concern is that $c(t)$ merely tracks the per-position next-token Shannon entropy $H_t$ of the model. On a control panel of $120$ random chr22 windows ($719{,}000$ analysed positions) we compute $H_t$ from the post-norm logits and find an overall Spearman $\rho(c, H_t)\!=\!-0.079$ (per-context $|\rho|\!\le\!0.16$ except $5'$~UTR; \S\ref{sec:discussion}; full per-context numbers in Tab.~\ref{tab:entropy-decoup}, App.~\ref{app:entropy-decoup}). On the same chr22 panel the splice-donor-versus-intron Cohen's $d$ is $-0.452$; after regressing $c$ on $H_t$ and re-measuring on the residual it strengthens to $-0.583$. Next-token entropy therefore does not explain the splice depth ordering.

\subsection{Motif Edits Separate Detection from Flanking Context}
\label{sec:motif-controls}

We probe $c(t)$ with two complementary perturbations on the same $1{,}000$ canonical GT-AG donors (chr22, $\pm 10$~bp pad-averaged $\overline{c}$). Real donors settle at $\overline{c}\!=\!26.77$ (Table~\ref{tab:motif-flank-main}). Replacing the central GT with AA while preserving the $\pm 100$~bp flank deepens $\overline{c}$ by $0.46$ layers (paired Wilcoxon $p\!=\!2.3\!\times\!10^{-32}$, $|d|\!=\!0.086$): a small but reliable signal that the model uses the canonical dinucleotide for early commitment. The GT$\!\to\!$AA edit also changes composition; a composition-preserving variant (e.g.\ GT$\!\to\!$TG) would isolate motif identity from composition. Dinucleotide-shuffling the $\pm 100$~bp flank while preserving the central GT produces the opposite shift, lifting $\overline{c}$ to a shallower value by $3.18$ layers (paired $p\!=\!4.1\!\times\!10^{-59}$, $|d|\!=\!0.515$): the isolated GT becomes easier to stabilise, whereas the real donor context requires deeper flanking-grammar integration. The two perturbations move $\overline{c}$ in opposite directions on the same donor loci, separating motif detection from context integration along the same axis. We cannot rule out from depth alone an alternative reading under which the model \emph{disengages} from uninformative noise on shuffled flanks and commits early to the lone GT (``early simplification'' rather than reduced integration); the within-splice motif breakdown in App.~\ref{app:splice-canonical} carries the same cautionary logic.

\begin{table}[!htbp]
  \caption{Motif-edit and flank-shuffle controls on $1{,}000$ canonical GT-AG donors (chr22, $\pm 10$\,bp pad-averaged $\overline{c}$). Paired Cohen's $|d|$ is $0.086$ for the GT motif edit and $0.515$ for the flank shuffle.}
  \label{tab:motif-flank-main}
  \centering
  \small
  \begin{tabular}{@{}lccc@{}}
    \toprule
    Condition & $\overline{c}$ & $\Delta \overline{c}$ & paired $p$ \\
    \midrule
    real GT-AG donor              & $26.77$ & $0$              & --- \\
    GT replaced with AA           & $27.24$ & $+0.46$          & $2.3\!\times\!10^{-32}$ \\
    flank shuffled, GT preserved  & $23.59$ & $\mathbf{-3.18}$ & $4.1\!\times\!10^{-59}$ \\
    \bottomrule
  \end{tabular}
\end{table}

\subsection{Variant-Disruption Depth Correlates with Molecular Consequence}
\label{sec:variant-disruption}

For variants we read the same cosine trajectory differentially as $\DDcos(\ell,t)\!=\!\Dcos^{\mathrm{alt}}(\ell,t) - \Dcos^{\mathrm{ref}}(\ell,t)$ and record the layer at which $|\DDcos|$ peaks. This converts the \gDTR{} lens from an absolute settling readout into a sensitivity readout. Across $4{,}023$ ClinVar P/LP variants in 15 cancer-associated genes, molecular-consequence classes show different peak-disruption distributions, with class medians ordered \emph{intron $<$ frameshift $<$ nonsense $<$ missense $\approx$ canonical splice $<$ synonymous} (Fig.~\ref{fig:variants}). The Kruskal--Wallis omnibus is significant ($p\!=\!3.0\!\times\!10^{-10}$, $H\!=\!53.2$)~\cite{kruskal1952wallis}, but the effect is small ($\varepsilon^2\!=\!0.013$) and the class distributions overlap broadly; only the nonsense-versus-missense contrast is individually resolved after Bonferroni-corrected Dunn testing ($p_{\mathrm{adj}}\!<\!10^{-3}$)~\cite{dunn1964multiple}. Pathogenic synonymous variants often act through splicing, so the synonymous tail may carry the canonical-splice signal rather than a separate synonymous one. We read this as a population-level depth shift rather than class separation or a pathogenicity score. Full $\DDcos$ trajectory AUROC analyses, cohort-construction details, and rank-biserial effect sizes are in App.~\ref{app:auroc-detail}.

\begin{figure}[!t]
  \centering
  \includegraphics[width=0.98\columnwidth]{figures/fig_3_variants.png}
  \caption{Variant consequences peak at different disruption layers. Argmax layer of $|\DDcos|$ across $4{,}023$ ClinVar P/LP variants in 15 cancer-associated genes, grouped by molecular consequence; class-median layer in red. Kruskal--Wallis 6-way $p\!=\!3.0\!\times\!10^{-10}$; the largest adjacent jump is between nonsense and missense ($p_{\mathrm{adj}}\!<\!10^{-3}$). The broad overlap reflects a population-level depth shift rather than class separation.}
  \label{fig:variants}
\end{figure}

\subsection{Robustness and Tokenisation Limits}
\label{sec:cross-arch}

Applying the same chr22 protocol to Evo~2~7B, HyenaDNA-large, NT-v2~500M, and DNABERT-2 gives a bounded robustness result (Table~\ref{tab:cross-arch-summary}; details in App.~\ref{app:cross-arch}). All four baselines use official HuggingFace checkpoints with frozen weights and the same regional $q_{70}$ calibration protocol. Donor settles earlier than intron in both per-base causal LMs, and their per-window settling depths are moderately concordant across models (Spearman $\rho\!=\!+0.516$). The MLMs operate at $k$-mer or BPE token granularity and cannot resolve single-base splice junctions. The cross-model result supports \gDTR{} within per-position causal LMs while identifying tokenisation granularity as the current architectural limit.

\begin{table}[t]
  \caption{Cross-model summary. Full per-window numbers and MLM tokenisation notes are in App.~\ref{app:cross-arch}.}
  \label{tab:cross-arch-summary}
  \centering
  \footnotesize
  \setlength{\tabcolsep}{4pt}
  \resizebox{\columnwidth}{!}{%
    \begin{tabular}{@{}lcc@{}}
      \toprule
      Family (model) & Per-pos.\ readout & Donor$\,{<}\,$intron \\
      \midrule
      Causal LM \scriptsize(Evo~2, HyenaDNA-large) & yes & yes \\
      MLM \scriptsize(NT-v2, DNABERT-2)            & no (per-window) & resolution-limited \\
      \bottomrule
    \end{tabular}
  }
\end{table}

%%%%%%%%%%%%%%%%%%%%%%%%%%%%%%%%
% 4. DISCUSSION
%%%%%%%%%%%%%%%%%%%%%%%%%%%%%%%%
\section{Discussion and Limitations}
\label{sec:discussion}

Settling depth is bidirectional, and this has a cost as well as a use. The use is that it explains within-class orderings that run against naive canonicity priors (Apps.~\ref{app:splice-canonical}, \ref{app:promoter-elements}) and warns against reading $c(t)$ as a scale of biological signal strength. The cost is that a single $c(t)$ does not interpret itself: a shallow settling can mean early motif commitment, a highly determinate context, or disengagement from uninformative input. The two matched perturbations of \S\ref{sec:motif-controls} move $c$ in opposite directions: disrupting the central GT deepens $c$, whereas shuffling the flanking grammar makes the preserved GT settle earlier. The context arm moves $c$ about six times as far as the motif arm ($|d|\!=\!0.515$ vs $0.086$), so on those loci $c(t)$ behaves mainly as a contextual-determinacy readout; whether the same balance holds at other functional contexts is open.

The remaining limitations group into four kinds. The lens is direction-only and discards residual magnitude (App.~\ref{app:lim-magnitude}), and the absolute readout saturates past the rotation layer, so settling and variant scoring are not read from the same layers (\S\ref{sec:variant-disruption}, App.~\ref{app:auroc-detail}). The motif and flank edits of \S\ref{sec:motif-controls} are off-manifold and cannot, on their own, separate reduced integration from early simplification; the within-splice ordering in App.~\ref{app:splice-canonical} is the on-manifold version of the same test. Composition and $k$-mer rarity are only partly controlled (App.~\ref{app:lim-composition}): the splice effect survives entropy and GC matching, but promoter and flank-shuffle contrasts remain composition-entangled (App.~\ref{app:promoter-elements}). Scope is narrow: a single calibration chromosome transferred to one held-out chromosome, replication on the only other per-base model available (HyenaDNA; Caduceus~\cite{schiff2024caduceus} is the obvious next per-base architecture to test), and a variant cohort of 15 cancer-associated genes.

The variant-depth ordering hints at an analogy with hierarchical encoding in protein language models~\cite{banerjee2025neuronlabelling}, but encoding depth, settling depth, and sensitivity-peak depth are distinct quantities, so we treat this as a hypothesis rather than a result. A concrete next step is to test it against public Evo~2 sparse autoencoders~\cite{deng2025interpreting}: do early-settling positions align with the activations of specific splice or exon SAE features? The rotation boundary itself is a portable methodological result. A residual-stream lens taken from NLP needs both a per-model reference and an explicit choice between absolute and differential readouts on genomic models (\S\ref{sec:framework}, App.~\ref{app:auroc-detail}), alongside causal or sparse-feature interventions~\cite{cunningham2023sparse,meng2022rome} that connect depth signatures to model circuits.

%%%%%%%%%%%%%%%%%%%%%%%%%%%%%%%%
% 5. CONCLUSION
%%%%%%%%%%%%%%%%%%%%%%%%%%%%%%%%
\section{Conclusion}
\label{sec:conclusion}

\gDTR{} provides a training-free, layer-resolved probe of where genomic foundation models stabilise token representations relative to their output-ready residual frame. Across Evo~2 analyses, settling depth is chromosome-transferable, shifts earlier at splice sites and enhancer-like cCREs, and responds in opposite directions to motif edits and flank shuffles, indicating that depth reflects both motif detection and contextual integration. Differential \gDTR{} further shows that variant-induced residual disruptions peak at consequence-associated layers, suggesting a complementary axis for interpreting variant effects rather than a standalone pathogenicity score. Future work should extend the analysis to whole-genome matched controls, additional per-base architectures, and causal or sparse-feature interventions~\cite{cunningham2023sparse,meng2022rome} that connect depth signatures to concrete model circuits.

%%%%%%%%%%%%%%%%%%%%%%%%%%%%%%%%
% BIBLIOGRAPHY
%%%%%%%%%%%%%%%%%%%%%%%%%%%%%%%%
\bibliography{gdtr_paper}
\bibliographystyle{icml2026}

%%%%%%%%%%%%%%%%%%%%%%%%%%%%%%%%
% APPENDIX
%%%%%%%%%%%%%%%%%%%%%%%%%%%%%%%%
\newpage
\appendix
\onecolumn

% Re-number appendix figures and tables as A1, A2, ... (independent
% of the main-text counters).
\renewcommand{\thefigure}{A\arabic{figure}}
\renewcommand{\thetable}{A\arabic{table}}
\setcounter{figure}{0}
\setcounter{table}{0}

\noindent{\LARGE\bfseries Appendix\par}
\vspace{0.6em}
\noindent{\large\itshape Supplementary material for ``\gDTR{}: Layer-wise Settling Depth Reveals Biological Grammar in Genomic Foundation Models''\par}
\vspace{1em}

%%%%%%%%%%%%%%%%%%%%%%%%%%%%%%%%
% APPENDIX A — Method details
%%%%%%%%%%%%%%%%%%%%%%%%%%%%%%%%
\section{Method Details}
\label{app:method}

\subsection{Addressing the Structural Characteristics of Evo~2: A Non-Functional Final Layer}
\label{app:architectural-quirk}

We observe that the final attention layer of Evo~2~7B is structurally non-functional with respect to the residual stream. Direct measurement on chr22 control sequences (100 windows of 6\,kb each) yields the values in Table~\ref{tab:idle-block}. Two observations are relevant: (i) $\max_t |h_{30}(t) - h_{31}(t)| = 0$ exactly, so layer~31 acts as a residual passthrough; (ii) the cosine alignment with the post-final-norm reference $\hnorm$ shifts from $\cos(h_{29}, \hnorm) = -0.013$ (near-orthogonal) to $\cos(h_{30}, \hnorm) = \cos(h_{31}, \hnorm) = +0.6855$, indicating that layer~30 performs the rotation into the post-norm direction. Consequently, the deepest tap not yet aligned with the output-ready frame is $L^{\star} = 29$. The standard NLP convention of tapping the final layer therefore does not apply to Evo~2; we instead use the post-final-RMSNorm state $\hnorm$ as the convergence reference and report $L^{\star} = 29$ as the empirical boundary of the pre-rotation taps.

\paragraph{Why the rotation matters as a measurement boundary.} This boundary has two consequences for the framework. First, the standard NLP recipe of reading the final layer would compare against a passthrough and an already-saturated state, which collapses the signal, so the post-final-norm reference is needed rather than cosmetic. Second, it places the settling readout in the pre-rotation taps ($\ell \le 29$). Once layer~30 aligns the stream with $\hnorm$, the absolute cosine distance saturates and no longer separates tokens. That is a property of the absolute readout, not evidence that later layers are inert: read \emph{differentially}, layer~30 carries the sharpest variant signal (\S\ref{sec:variant-disruption}, App.~\ref{app:auroc-detail}). The rotation is therefore a boundary in measurement geometry, with level readouts useful before it and sensitivity readouts after it. It is also the main reason a residual-stream lens taken from NLP needs a per-model reference on genomic models.

\begin{table}[!htbp]
  \caption{Evidence that the final attention layer of Evo~2 is structurally non-functional with respect to the residual stream. Layer~31 is bit-identical to layer~30 in value but stored as a distinct tensor; both align cosine-wise with the post-final-RMSNorm reference, identifying layer~29 as the deepest tap not yet aligned with the output frame.}
  \label{tab:idle-block}
  \centering
  \small
  \begin{tabular}{@{}lcc@{}}
    \toprule
    Quantity & Measured value & Interpretation \\
    \midrule
    $\max\bigl|h_{30}-h_{31}\bigr|$ & $0$ (exact) &
    layer 31 is residual passthrough \\
    \texttt{data\_ptr}($h_{30}$) vs.\ $h_{31}$ & distinct &
    physically separate copies \\
    $\cos(h_{31}, \hnorm)$ & $0.6855$ &
    post-norm differs from raw $h_{31}$ \\
    $\cos(h_{30}, \hnorm)$ & $0.6855$ &
    identical to $h_{31}$ \\
    $\cos(h_{29}, \hnorm)$ & $-0.013$ &
    deepest distinct tap ($L^{\star}=29$) \\
    \bottomrule
  \end{tabular}
\end{table}

\subsection{Hyperparameter Sensitivity}
\label{app:gamma-sweep}

On chr22 the regional $q_{70}$ calibration yields $\gamma_{\cos} = 0.397$. A $5 \times 5$ grid search over $(\gamma_{\cos}, \rho)$ (Table~\ref{tab:gamma-sweep}) produces a wide, flat plateau around the operating point $(\gamma_{\cos} = 0.40, \rho = 0.85)$ with $\pm 0.10 \times \pm 0.05$ variation. Within this plateau the standard
deviation across the 25 cells is $0.06$, indicating low sensitivity to
small perturbations in either hyperparameter. This robustness was first
identified in Phase~0 (HyenaDNA-medium-160k), where the analogous
$q_{70}$ calibration produced a best operating point of
$(\gamma_{\cos}\!=\!0.50,\; \rho\!=\!0.85)$ with Cohen's $d = -1.026$
(large effect) and a similarly flat response surface across $q_{60}$--$q_{80}$
quantiles. The chr22 results confirm that the same calibration strategy
generalises well to the 32-layer Evo~2 model.

\begin{table}[!htbp]
  \caption{Effect-size grid search over $(\gamma_{\cos}, \rho)$. Values are reported on the analysis scale used for hyperparameter calibration. The locked
    operating point $(0.40, 0.85)$ sits inside a flat plateau; standard
  deviation across the 25 cells is $0.06$.}
  \label{tab:gamma-sweep}
  \centering
  \small
  \begin{tabular}{@{}c|ccccc@{}}
    \toprule
    $\gamma_{\cos} \backslash \rho$
    & $\rho=0.70$ & $\rho=0.75$ & $\rho=0.80$ & $\rho=0.85$ & $\rho=0.90$ \\
    \midrule
    $0.30$ & $5.04$ & $5.18$ & $5.21$ & $5.20$ & $5.15$ \\
    $0.35$ & $5.10$ & $5.21$ & $5.24$ & $5.23$ & $5.18$ \\
    $0.40$ & $5.16$ & $5.25$ & $\mathbf{5.28}$ & $\mathbf{5.26}$ & $5.21$ \\
    $0.45$ & $5.13$ & $5.22$ & $5.25$ & $5.24$ & $5.19$ \\
    $0.50$ & $5.08$ & $5.17$ & $5.20$ & $5.19$ & $5.14$ \\
    \bottomrule
  \end{tabular}
\end{table}

\subsection{Entropy Decoupling per Context}
\label{app:entropy-decoup}

Table~\ref{tab:entropy-decoup} reports per-context Spearman $\rho(c,
H_t)$ on a $120$-window chr22 control panel ($719{,}000$ analysed positions; the
  $1{,}000$ shortfall versus $120\!\times\!6{,}000$ reflects truncated logits at
window edges where the causal context is incomplete), where
$H_t$ is the per-position next-token Shannon entropy from the post-norm
logits. The overall correlation is small and slightly negative
($\rho\!=\!-0.079$); per-context correlations stay $|\rho|\!\leq\!0.16$
for every context except 5$'$~UTR ($\rho\!=\!+0.41$). The largest
non-UTR couplings are splice acceptor at $|\rho|\!=\!0.152$ and intron
at $|\rho|\!=\!0.108$ (Tab.~\ref{tab:entropy-decoup}). After regressing
$c$ on $H_t$ on the same chr22 panel, the splice-donor-versus-intron
Cohen's $d$ \emph{strengthens} from $-0.452$ to $-0.583$ on the
residual, confirming that next-token entropy does not explain the
splice depth ordering.

\begin{table}[!htbp]
  \caption{Per-context Spearman correlation between settling depth $c(t)$
    and next-token entropy $H_t$ on the chr22 control panel
    ($120$ windows, $719{,}000$ analysed positions; the $1{,}000$ shortfall versus
    $120\!\times\!6{,}000$ reflects truncated logits at window edges). The next-largest
    non-UTR couplings are splice acceptor ($|\rho|\!=\!0.152$) and intron
  ($|\rho|\!=\!0.108$); 5$'$~UTR is the only entropy-coupled context.}
  \label{tab:entropy-decoup}
  \centering
  \small
  \begin{tabular}{@{}lcccc@{}}
    \toprule
    Context & $n$ & $\overline{c}$ & $\bar H_t$ & $\rho(c, H_t)$ \\
    \midrule
    intergenic        & $243{,}956$ & $28.94$ & $0.795$ & $-0.024$ \\
    intron            & $425{,}131$ & $28.03$ & $1.030$ & $-0.108$ \\
    coding exon       & $\phantom{0}36{,}427$ & $28.61$ & $0.744$ & $-0.080$ \\
    3$'$ UTR          & $\phantom{00}8{,}238$ & $27.61$ & $1.073$ & $-0.028$ \\
    splice donor      & $\phantom{00}1{,}981$ & $25.04$ & $0.887$ & $-0.083$ \\
    splice acceptor   & $\phantom{00}1{,}920$ & $27.33$ & $0.764$ & $-0.152$ \\
    \textbf{5$'$ UTR} & $\phantom{00}2{,}347$ & $31.43$ & $1.121$ & $\mathbf{+0.414}$ \\
    \midrule
    overall           & $719{,}000$ & --- & --- & $-0.079$ \\
    \bottomrule
  \end{tabular}
\end{table}

\subsection{Tuned-Lens Recovery Across All 32 Layers}
\label{app:tuned-lens}

Each layer is fitted with a single $4096 \times 4096$ affine $A_{\ell}$
using MSE between $A_{\ell} h_{\ell}$ and $\hnorm$, optimised with Adam
($\beta_1 = 0.9$, $\beta_2 = 0.999$, weight decay $= 0$) at learning
rate $10^{-3}$ for 15 epochs over 100 calibration sequences with batch
size 1. The affine weights are used only for this verification and are
discarded at inference, so no hyperparameter governs the production
settling-depth pipeline. The
post-norm tap is the prediction target. Recovery scores at five
representative layers are reported in Table~\ref{tab:tuned-lens};
across the full 32-layer distribution, $30$ reach $\ge 98\%$, with the
two layers below this threshold not among the five shown (so the
table's ``worst of these five'' label refers to the displayed subset,
not the overall minimum). The $98\%$ threshold is descriptive of
the empirical recovery distribution rather than a pre-registered
cut-off: it summarises where the bulk of the per-layer recovery scores
sit, not a criterion that the framework is required to clear at
inference (the framework is training-free and uses no affine weights
when computing $c(t)$).

\begin{table}[!htbp]
  \caption{Tuned-lens recovery at five representative layers spanning
  network depth. Block-type abbreviations follow Evo~2's StripedHyena~2
  architecture: \texttt{hcl} (Hyena-Convolution-Long), \texttt{hcm}
  (Hyena-Convolution-Medium), \texttt{hcs} (Hyena-Convolution-Short).}
  \label{tab:tuned-lens}
  \centering
  \small
  \begin{tabular}{@{}lccc@{}}
    \toprule
    Layer / block type & Initial MSE & Final MSE (15 epochs) & Recovery \\
    \midrule
    $\ell=2$ (\texttt{hcl}) & $1{,}259$ & $5.6$ & $0.9956$ \\
    $\ell=12$ (\texttt{hcm}) & $742$ & $13.6$ & $0.9816$ (worst of these five) \\
    $\ell=22$ (\texttt{hcm}) & $418$ & $0.41$ & $0.9990$ \\
    $\ell=28$ (\texttt{hcs}) & $822$ & $0.34$ & $0.9996$ (best below tap) \\
    $\ell=29$ (canonical tap, \texttt{hcm}) & $510$ & $0.20$ & $0.9996$ \\
    \bottomrule
  \end{tabular}
\end{table}

%%%%%%%%%%%%%%%%%%%%%%%%%%%%%%%%
% APPENDIX B — Cross-architecture replication (merged: model specs,
% per-context tables, two-tier figure)
%%%%%%%%%%%%%%%%%%%%%%%%%%%%%%%%
\section{Cross-Architecture Replication: Scope, Granularity, Two-Tier Structure}
\label{app:cross-arch}

This appendix expands the bounded robustness result of
\S\ref{sec:cross-arch}. We compare four publicly released genomic
foundation models that span the per-bp causal-LM and the tokenised MLM
design space: Evo~2~7B, HyenaDNA-large-1m, NT-v2~500M, and
DNABERT-2~117M. Architectural specifications, the per-window token budget
they impose on the same chr22 sequences, end-to-end runtime, and the
per-model $q_{70}$-calibrated $\gamma_{\cos}$ are summarised in
Table~\ref{tab:fm-specs}. All four models use the identical chr22
$12{,}978$-window panel.

\begin{table}[!htbp]
  \caption{Architectural specifications, tokenisation, per-window context,
    runtime, and per-model $q_{70}$-calibrated $\gamma_{\cos}$ for the four
  genomic foundation models compared in \S\ref{sec:cross-arch}.}
  \label{tab:fm-specs}
  \centering
  \footnotesize
  \setlength{\tabcolsep}{4pt}
  \begin{tabular}{@{}llcccll@{}}
    \toprule
    Model & Architecture & Layers & Hidden & Tokens / 6\,kb window
    & Wall-clock & $\gamma_{q_{70}}$ \\
    \midrule
    Evo~2~7B          & Hybrid Transformer + StripedHyena~2 & $32$ & $4096$
    & $6{,}000$ (1\,bp)         & reused      & $0.397$ \\
    HyenaDNA-large-1m & Pure Hyena                          & $\phantom{0}8$ & $\phantom{0}256$
    & $6{,}001$ (1\,bp + BOS)   & $\approx 4$\,min    & $0.358$ \\
    NT-v2~500M        & Transformer MLM ($k$-mer)           & $29$ & $1024$
    & $\phantom{0,}671$ ($k=6$, $\approx 4$\,kb) & $\approx 7.5$\,min  & $0.533$ \\
    DNABERT-2~117M    & Transformer MLM (BPE)               & $12$ & $\phantom{0}768$
    & $\approx 600$ (BPE, $\approx 3$\,kb)        & $\approx 3$\,min    & $0.677$ \\
    \bottomrule
  \end{tabular}
\end{table}

\paragraph{Per-context replication.} Table~\ref{tab:cross-arch-context}
reports the per-context mean settling depth for each model. The two
per-bp causal LMs (Evo~2, HyenaDNA-large) replicate
donor settles earlier than intron, and acceptor earlier than intron, at depth-normalised scale
(approximately 3 layers below intron in 32-layer Evo~2; approximately 0.34 layers below intron in 8-layer HyenaDNA). The two MLMs tokenise at $k$-mer/BPE
granularity: their per-window readout collapses splice-junction signal
into the surrounding window mean and therefore matches the exon-dominant
window mean to within rounding, consistent with a tokenisation-induced
loss of single-base resolution rather than a model-class disagreement.

\begin{table}[!htbp]
  \caption{Per-context mean settling depth across the four models on the
    identical chr22 $12{,}978$-window set. Causal-LM models recover
    donor/acceptor~$<$~intron at depth-normalised scale; MLM models
  report per-window estimates with no per-position separation.}
  \label{tab:cross-arch-context}
  \centering
  \footnotesize
  \setlength{\tabcolsep}{4pt}
  \begin{tabular}{@{}llccccc@{}}
    \toprule
    Model & Kind & $L$ & $\gamma_{q_{70}}$
    & $\overline{c}$ (intron) & $\overline{c}$ (coding exon)
    & $\overline{c}$ (donor / acceptor) \\
    \midrule
    Evo~2~7B            & per\_position & $32$ & $0.397$
    & $27.84$ & $28.28$ & $25.59 \,/\, 25.71$ \\
    HyenaDNA-large      & per\_position & $8$  & $0.358$
    & $\phantom{0}6.89$ & $\phantom{0}6.67$
    & $\phantom{0}6.55 \,/\, \phantom{0}6.62$ \\
    NT-v2~500M          & per\_window   & $29$ & $0.533$
    & n.a. (intron-dominant 0)
    & $27.80$ (exon-dom.) & $27.85$ (splice-cont.) \\
    DNABERT-2~117M      & per\_window   & $12$ & $0.677$
    & n.a. (intron-dominant 0)
    & $11.28$ (exon-dom.) & $11.27$ (splice-cont.) \\
    \bottomrule
  \end{tabular}
\end{table}

\paragraph{Pairwise rank concordance.} Table~\ref{tab:cross-arch-rho}
shows the pairwise Spearman $\rho$ of per-window mean settling depth
across the four models. The pattern is two-tier: within-family
correlations are positive ($\rho_{\mathrm{Evo2,Hyena}}\!=\!+0.516$;
$\rho_{\mathrm{NT,DNABERT}}\!=\!+0.663$), whereas every cross-family
pair is weakly negative. This is the empirical basis for the
``replication within per-bp causal LMs, tokenisation-bound for MLMs''
claim of \S\ref{sec:cross-arch}.

\begin{table}[!htbp]
  \caption{Pairwise Spearman $\rho$ of per-window mean settling depth
    across the four genomic foundation models (chr22 windows).
    Within-family correlations are positive; cross-family correlations are
  weakly negative.}
  \label{tab:cross-arch-rho}
  \centering
  \footnotesize
  \setlength{\tabcolsep}{4pt}
  \begin{tabular}{@{}lcccc@{}}
    \toprule
    & Evo~2~7B & HyenaDNA-large & NT-v2~500M & DNABERT-2 \\
    \midrule
    Evo~2~7B         & $1.000$  & $+0.516$ & $-0.119$ & $-0.188$ \\
    HyenaDNA-large   & $+0.516$ & $1.000$  & $-0.287$ & $-0.166$ \\
    NT-v2~500M       & $-0.119$ & $-0.287$ & $1.000$  & $+0.663$ \\
    DNABERT-2        & $-0.188$ & $-0.166$ & $+0.663$ & $1.000$  \\
    \bottomrule
  \end{tabular}
\end{table}

\paragraph{Visual summaries.} The per-context bar chart
(Fig.~\ref{fig:cross-arch-context}) shows the
donor/acceptor~$<$~intron inequality directly for the per-bp causal
LMs and the absence of separation for the tokenised MLMs.
Fig.~\ref{fig:cross-arch} consolidates the rank-concordance heatmap, the
depth-normalised donor-vs-intron comparison, and a schematic of the
two-tier structure into a single panel.

\begin{figure}[!htbp]
  \centering
  \includegraphics[width=0.95\linewidth]{figures/fig_A1_cross_arch_context.png}
  \caption{Per-context mean settling depth across the four genomic
    foundation models. Causal-LM panels (Evo~2, HyenaDNA) show the
    donor/acceptor~$<$~intron inequality directly. MLM panels (NT-v2,
  DNABERT-2) show no separation at tokenised window resolution.}
  \label{fig:cross-arch-context}
\end{figure}

\begin{figure}[!htbp]
  \centering
  \includegraphics[width=0.95\linewidth]{figures/fig_A2_cross_arch_two_tier.png}
  \caption{Cross-architecture two-tier structure (numerical detail in
    Tables~\ref{tab:cross-arch-context}, \ref{tab:cross-arch-rho}).
    \textbf{(a)}~Pairwise Spearman $\rho$ of per-window mean settling depth
    across four genomic foundation models on chr22, showing the within-family
    positive block-diagonal vs.\ cross-family weak-negative off-diagonal
    structure. \textbf{(b)}~Schematic of the two-tier structure with the
    within/cross-family $\rho$ values, summarising the
    ``per-bp causal LMs replicate, tokenised MLMs are bound by readout
  granularity'' interpretation.}
  \label{fig:cross-arch}
\end{figure}

%%%%%%%%%%%%%%%%%%%%%%%%%%%%%%%%
% APPENDIX C — Variant pathogenicity AUROC (merged: per-layer ablation
% + AUROC summary + four-panel figure)
%%%%%%%%%%%%%%%%%%%%%%%%%%%%%%%%
\FloatBarrier
\section{Variant Pathogenicity: AUROC Summary, Per-Layer Ablation, and Panels}
\label{app:auroc-detail}

This appendix supports the trajectory-information verification in \S\ref{sec:variant-disruption}. The cohort is the $8{,}008$ ClinVar P/LP-vs-B/LB SNVs across 15 cancer-associated genes; the splitter is 10-fold stratified cross-validation with seed $42$, and the LOGO-CV column reports leave-one-gene-out generalisation across the 14 evaluable genes. We treat AUROC strictly as an empirical validation that the 32-d $\DDcos$ trajectory carries information beyond any single tap; we do not position \gDTR{} as a clinical scorer (\S\ref{sec:variant-disruption}).

\begin{table}[!htbp]
  \caption{Feature-level AUROC summary. The primary cosine lens reaches
    $0.844$ with the full 32-d trajectory and only $0.729$ at its best single
    tap -- a $+0.115$ gap. Bracketed numbers are $1{,}000$-bootstrap $95\%$
  confidence intervals.}
  \label{tab:auroc}
  \centering
  \small
  \begin{tabular}{@{}lccc@{}}
    \toprule
    Feature & dim. & Stratified 10-fold AUROC & LOGO-CV AUROC \\
    \midrule
    Best single-layer $\DDcos$ ($\ell=30$ tap) & $1$
    & $0.729$ $[0.717,\, 0.741]$
    & $0.726$ $[0.694,\, 0.758]$ \\
    Best single-layer $\DDjsd$ ($\ell=29$ canonical tap) & $1$
    & $0.794$ $[0.781,\, 0.806]$
    & $0.787$ $[0.752,\, 0.821]$ \\
    Evo~2 log-likelihood & $1$
    & $0.751$ $[0.738,\, 0.764]$
    & $0.793$ $[0.740,\, 0.846]$ \\
    32-d $\DDjsd$ vector & $32$
    & $0.823$ $[0.813,\, 0.832]$
    & $0.821$ $[0.790,\, 0.853]$ \\
    \textbf{32-d $\DDcos$ vector (primary)} & $32$
    & $\mathbf{0.844}$ $[0.831,\, 0.857]$
    & $\mathbf{0.843}$ $[0.811,\, 0.876]$ \\
    Ensemble ($\Delta D + $ Evo~2 LL) & $33$
    & $0.861$ $[0.851,\, 0.871]$
    & $0.866$ $[0.832,\, 0.899]$ \\
    \bottomrule
  \end{tabular}
\end{table}

\paragraph{Per-layer ablation.} Table~\ref{tab:per-layer} reports
single-layer logistic-regression AUROC at representative taps for both
lenses. The two lenses agree on the U-shape of layer-wise discriminative
mass but disagree on which tap is sharpest: JSD concentrates mass at the
canonical tap $\ell\!=\!29$, whereas cosine spreads mass across many
taps and accordingly produces a much larger 32-d-vector-vs-best-single-tap
gap.

\begin{table}[!htbp]
  \caption{Selected per-layer AUROCs. JSD concentrates discriminative
    mass at the canonical tap $\ell\!=\!29$; cosine spreads it across many
  taps, producing a much larger 32-d-vs-best-single-tap gap.}
  \label{tab:per-layer}
  \centering
  \small
  \begin{tabular}{@{}llcc@{}}
    \toprule
    Layer & Block type & $\DDjsd$ AUROC & $\DDcos$ AUROC \\
    \midrule
    $0$  & embed                    & $0.605$ & $0.519$ \\
    $7$  & \texttt{hcs} (shallow)   & $0.662$ & $0.595$ \\
    $12$ & \texttt{hcm}             & $0.656$ & $0.555$ \\
    $17$ & \texttt{attn}            & $0.723$ & $0.646$ \\
    $24$ & \texttt{attn}            & $0.685$ & $0.612$ \\
    $28$ & \texttt{hcs}             & $0.565$ & $0.698$ \\
    $29$ (canonical tap) & \texttt{hcm}
    & $\mathbf{0.794}$ (best JSD) & $0.604$ \\
    $30$ (post-norm$-1$)  & ---     & $0.512$ & $\mathbf{0.729}$ (best cos) \\
    $31$ (non-functional, see App.~\ref{app:architectural-quirk}) & ---
    & $0.499$ (degen.) & $0.729$ ($=h_{30}$) \\
    \midrule
    32-d vector & all                 & $0.823$ & $0.844$ \\
    \bottomrule
  \end{tabular}
\end{table}

\paragraph{Two readings of one lens: why settling and variant scoring favour different taps.}
Settling depth and variant scoring read the same cosine lens in two
different ways, and this is what makes them favour different layers.
Settling reads the \emph{absolute} distance $\Dcos(\ell,t)$ as a level
crossing: it is degenerate at any tap whose representation already
matches the reference, so once layer~30 rotates the residual stream into
$\hnorm$ the absolute distance collapses to a narrow floor
($\Dcos(30,\hnorm)\!\approx\!0.31$, not $0$) and can no longer separate
tokens. Variant scoring instead reads the \emph{difference}
$\DDcos(\ell,t)$ as a sensitivity: subtracting the shared post-rotation
floor is precisely what isolates the allele-specific residual, so the
signal is sharpest exactly where the absolute distance is least
informative. The two readings therefore predict opposite
layer-optimality, which is why the single-layer cosine AUROC is $0.604$
at $L^{\star}\!=\!29$ but $0.729$ at $\ell\!=\!30$: the gap reflects one
quantity being read before the rotation step and the other after it, not
that $L^{\star}\!=\!29$ is a worse tap. The reference here is $\hnorm$, the
post-final-RMSNorm state, not $h_{30}$ itself; what saturates at
layer~30 is the across-token spread of the absolute distance, which is
harmless for the $\Delta$-based sensitivity but fatal for level-based
settling. The $\ell\!=\!31$ AUROC inherits the same $0.729$ because
layer~31 is a residual passthrough (Tab.~\ref{tab:idle-block}); it is
reported for completeness rather than as new information.

\paragraph{ROC, DeLong, and per-gene panels.}
Fig.~\ref{fig:auroc} provides the four diagnostic panels referenced
above. ROC curves (a) confirm the ranking from
Table~\ref{tab:auroc}; the Paired DeLong tests in (b) show that
$\DDcos$ adds significant information beyond Evo~2 $\Delta\textrm{LL}$
($p\!=\!3.6\!\times\!10^{-15}$); (c) is the per-layer ablation visualised
across all 32 taps, and (d) is the leave-one-gene-out AUROC across the
14 evaluable genes (uniformly $\ge 0.77$).

\begin{figure}[!htbp]
  \centering
  \includegraphics[width=0.95\linewidth]{figures/fig_A3_auroc.png}
  \caption{Variant pathogenicity discrimination on $8{,}008$ ClinVar
    P/LP vs B/LB SNVs across 15 cancer-associated genes (numerical summary
    in Table~\ref{tab:auroc}). \textbf{(a)}~ROC curves for the four scoring
    features. \textbf{(b)}~Paired DeLong tests: $\Delta D$ adds
    significant information beyond Evo~2 $\Delta\textrm{LL}$
    ($p\!=\!3.6\!\times\!10^{-15}$). \textbf{(c)}~Per-layer single-feature
    AUROC across all 32 layers -- best single tap is $\ell\!=\!29$ for JSD
    and $\ell\!=\!30$ for cosine (the latter is the post-norm rotation
      layer, not the canonical pre-rotation tap; see
    App.~\ref{app:architectural-quirk}); the 32-d vector beats either by
    $+0.05$ to $+0.12$.
    \textbf{(d)}~Leave-one-gene-out AUROC across 14 evaluable genes is
  uniformly high ($\ge 0.77$).}
  \label{fig:auroc}
\end{figure}

%%%%%%%%%%%%%%%%%%%%%%%%%%%%%%%%
% APPENDIX D — Splice anatomy beyond the headline contexts
% (merged: positional fine-profile + canonical-vs-non-canonical motifs)
%%%%%%%%%%%%%%%%%%%%%%%%%%%%%%%%
\section{Splice Anatomy Beyond the Headline Contexts}
\label{app:splice-anatomy}

This appendix supports the cautionary message in
\S\ref{sec:motif-controls}: the splice signal is real, but its internal
structure is asymmetric and motif-class-dependent.

\subsection{Positional Fine-Profile Around Donor / Acceptor}
\label{app:splice-fine}

On chr17 the donor profile reaches its mean settling-depth minimum at
position $+20$\,bp ($\overline{c}\!=\!24.06$) and remains $\ge 1.5$ layers
below the pooled intronic baseline ($\overline{c}\!=\!27.72$) out to $\pm 200$\,bp. On
the same chromosome the acceptor profile reaches its minimum at
$+50$\,bp ($\overline{c}\!=\!23.33$) and similarly remains depressed beyond
$\pm 200$\,bp. On chr22 both donor and acceptor minima sit at coordinate
$0$ ($\overline{c}\!=\!23.65$ and $23.64$ respectively). The asymmetry, with the donor minimum on the exonic side and the acceptor minimum on the intronic side, mirrors the asymmetric splice-grammar features (branch point and polypyrimidine tract) that lie predominantly on the intronic side of acceptors. The full
positional profile is shown in Fig.~\ref{fig:splice-fine}; the per-side
minima are summarised in Table~\ref{tab:splice-fine}.

\begin{figure}[!htbp]
  \centering
  \includegraphics[width=0.95\linewidth]{figures/fig_A4_splice_fine.png}
  \caption{Mean settling depth $\overline{c}$ as a function of distance to the
    nearest splice donor / acceptor, computed on the pooled chr17 + chr22
    analysis windows. Both profiles dip well below the intronic baseline
    ($\overline{c}\!=\!27.72$, red dashed) within $\pm 200$\,bp; donors minimise
    on the exonic side, acceptors on the intronic side, mirroring the
    asymmetric splice-grammar context (branch point and polypyrimidine
  tract on the intronic side of acceptors).}
  \label{fig:splice-fine}
\end{figure}

\begin{table}[!htbp]
  \caption{Per-side minima of the splice positional fine-profile.
    Pooled chr17 + chr22 analysis windows; positions measured from the
    splice junction. Donor exonic-side / acceptor intronic-side minima
  recapitulate the asymmetric splice-grammar context.}
  \label{tab:splice-fine}
  \centering
  \small
  \begin{tabular}{@{}lccccc@{}}
    \toprule
    Side & arg-min position & $\overline{c}$ at min & $\overline{c}$ at $-100$\,bp
    & $\overline{c}$ at $+100$\,bp & $n$ \\
    \midrule
    splice donor    & $+20$\,bp & $23.65$ & $26.05$ & $25.65$ & $280{,}944$ \\
    splice acceptor & $+50$\,bp & $23.64$ & $26.01$ & $24.50$ & $278{,}955$ \\
    \bottomrule
  \end{tabular}
\end{table}

\subsection{Canonical vs.\ Non-Canonical Splice Motif Breakdown}
\label{app:splice-canonical}

A finer dissection by motif class, using the genomic dinucleotides immediately downstream of the donor and upstream of the acceptor, reveals an unexpected ordering (Table~\ref{tab:splice-canonical}): \emph{non-canonical} splice donors converge \emph{earlier} than the dominant canonical GT-AG donors, while the minor canonical GC-AG class is the deepest of the three. The shallower-than-intron pattern
($\overline{c}\!=\!27.72$) holds for every splice class, but the within-class
ordering is the opposite of a naive ``stronger motif implies
shallower recognition'' prior. The Cohen's $d$ between canonical GT-AG
and non-canonical donors is small ($\approx 0.05$), consistent with the
constrained branch-point and polypyrimidine context that flanks
non-canonical introns~\cite{wang2008splicecode,burge1998u12}; we
report the magnitudes honestly rather than over-interpret.
% REVIEW FIX (concern #1, Table A11): turn the non-monotone ordering into
% evidence FOR the two-sided metric (arm (b), U12 tighter flank) rather than
% a contradiction; concedes "not monotone in signal strength" on our terms.
This ordering is a direct, positive
demonstration of arm (b) (context integration) acting independently of
motif canonicity. A metric that ranked motifs monotonically by signal strength would have predicted GT-AG $<$ GC-AG $<$ non-canonical in settling depth. The observed inversion, in which non-canonical donors (e.g.\ U12-type AT-AC introns) settle earliest because they sit in a tighter, more constrained flank that the model commits to early, is therefore not a contradiction but a clean within-splice instance of the bidirectional reading of \S\ref{sec:framework}. Settling depth tracks
contextual determinacy, not biological signal strength in any monotone
sense. The finding replicates qualitatively on the 8-layer HyenaDNA-large
at depth-normalised scale.

\begin{table}[!htbp]
  \caption{Within-splice motif breakdown by donor dinucleotide. The
    shallower-than-intron pattern ($\overline{c}\!=\!27.72$) holds for every
    splice class, but the within-class ordering is the opposite of a naive
    ``stronger motif implies shallower recognition'' prior:
  non-canonical donors converge earliest, GC-AG latest.}
  \label{tab:splice-canonical}
  \centering
  \small
  \begin{tabular}{@{}lccc@{}}
    \toprule
    Donor motif class & $n$ & $\overline{c}$ & vs.\ canonical GT-AG \\
    \midrule
    non-canonical (AT-AC, other) & $\phantom{00}5{,}977$ & $25.13$ & $d\approx -0.05$ (shallower) \\
    canonical GT-AG (dominant)   & $350{,}552$ & $25.79$ & --- (reference) \\
    canonical GC-AG (minor)      & $\phantom{00}5{,}165$ & $27.01$ & $d\approx +0.10$ (deeper) \\
    \midrule
    intron baseline              & $\approx\!10^{7}$\phantom{0} & $27.72$ & --- (deeper than every class) \\
    \bottomrule
  \end{tabular}
\end{table}

\subsection{Motif and Flank Perturbation Summary}
\label{app:motif-flank}

The motif-edit and flank-shuffle controls are reported in the body of the paper at Table~\ref{tab:motif-flank-main} alongside the detailed discussion in \S\ref{sec:motif-controls}. We retain this appendix anchor so that prior references to App.~\ref{app:motif-flank} resolve to the consolidated discussion.

%%%%%%%%%%%%%%%%%%%%%%%%%%%%%%%%
% APPENDIX E — Entropy decoupling at genome-wide scale
%%%%%%%%%%%%%%%%%%%%%%%%%%%%%%%%
\section{Entropy Decoupling at Genome-Wide Scale}
\label{app:restriction-of-range}

The body of the paper treats $c(t)$ as decoupled from per-position
next-token entropy $H_t$ at all canonical contexts except $5'$~UTR
(\S\ref{sec:shallowness}). Here we extend the diagnostic from the
$120$-window control panel of \S\ref{sec:shallowness} to the full
chr22 analysis windows ($12{,}978$ windows, $38.9$\,M analysed positions)
to confirm that the decoupling argument is not an artefact of the
small control sample. We use the canonical $c(t)$ from the chr22 forward
joined to a per-position $H_t$ computed from the post-norm logits on
the same chr22 sequence. Table~\ref{tab:restriction-of-range} reports
the per-context joint $(n, \overline{c}, \sigma(H_t), \mathrm{IQR}(H_t),
\rho(c, H_t))$.

Two points are worth noting. First, $|\rho(c, H_t)|\!\le\!0.224$ across
all seven canonical contexts at this scale; the entropy decoupling
argument therefore holds genome-wide. The largest residual couplings
appear at splice acceptor ($\rho\!=\!-0.224$) and intron
($\rho\!=\!-0.149$); these are negative, indicating that within these contexts higher prediction confidence is associated with shallower settling. This direction is consistent with the splice signal being driven by representational stability rather than by confidence inflation. Second,
$5'$~UTR yields $\rho\!=\!+0.054$ at $n\!=\!74{,}728$ positions, smaller
than the $+0.41$ observed on the $120$-window control sample. We
interpret the control-panel value as a small-sample property of the
specific window draw rather than a stable biological feature, and
continue to keep $5'$~UTR on a separate axis in Fig.~\ref{fig:shallowness}
as the most conservative reporting choice. A simple
$\sigma(H_t)$-driven restriction-of-range explanation (whereby narrow
entropy spread mechanically suppresses $\rho$) is not supported by the
data: $5'$~UTR has the narrowest $\sigma(H_t)$ in the table yet does
not produce the smallest $|\rho|$.

\begin{table}[!htbp]
  \caption{Per-context joint distribution of canonical settling depth
    $c(t)$ and per-position next-token entropy $H_t$ on the full chr22
    analysis windows ($38.9$\,M positions; central $3$\,kb of each
    $12{,}978$ window). $\sigma(H_t)$ and $\mathrm{IQR}(H_t)$ measure
    per-context entropy spread; $\rho(c, H_t)$ is Spearman correlation.
    Contexts are ordered by $|\rho|$.}
  \label{tab:restriction-of-range}
  \centering
  \small
  \begin{tabular}{@{}lccccc@{}}
    \toprule
    Context & $n$ & $\overline{c}$ & $\sigma(H_t)$ & $\mathrm{IQR}(H_t)$ & $\rho(c, H_t)$ \\
    \midrule
    splice acceptor   & $\phantom{0}{93{,}883}$ & $25.69$ & $0.496$ & $0.944$ & $-0.224$ \\
    intron            & $20{,}595{,}963$        & $27.82$ & $0.395$ & $0.563$ & $-0.149$ \\
    splice donor      & $\phantom{0}{94{,}528}$ & $25.59$ & $0.507$ & $0.982$ & $-0.134$ \\
    coding exon       & $\phantom{}1{,}643{,}533$ & $28.25$ & $0.504$ & $1.009$ & $-0.097$ \\
    5$'$~UTR          & $\phantom{0}{74{,}728}$ & $29.03$ & $0.366$ & $0.336$ & $+0.054$ \\
    3$'$~UTR          & $\phantom{00}{610{,}330}$ & $27.73$ & $0.357$ & $0.226$ & $-0.052$ \\
    intergenic        & $15{,}821{,}035$        & $28.75$ & $0.468$ & $0.918$ & $+0.001$ \\
    \bottomrule
  \end{tabular}
\end{table}

%%%%%%%%%%%%%%%%%%%%%%%%%%%%%%%%
% APPENDIX F — Single-base promoter elements
%%%%%%%%%%%%%%%%%%%%%%%%%%%%%%%%
\section{Single-Base Promoter Elements}
\label{app:promoter-elements}

To test whether the splice-site shallow-settling signature of
\S\ref{sec:shallowness} generalises to other single-base regulatory
anchors, we measure canonical $c(t)$ at four promoter-element classes
on chr22: canonical transcription start sites (TSS, GENCODE v44
Ensembl-canonical), and JASPAR position-weight-matrix hits for TATA-box
(MA0108.3, TBP), CAAT-box (MA0060, NFY), and GC-box (MA0079, Sp1). PWM
hits are restricted to positions within $\pm 1.5$\,kb of a canonical
TSS so that each anchor lies in a functional promoter context. Because
promoter elements are composition-loaded (GC-box is GC-rich by
definition), we report Cohen's $d$ against two baselines: the chr22
intron baseline used elsewhere in the paper, and a GC-content-matched
non-promoter chr22 control (5\%-wide GC bin, 1:1 matched).
Table~\ref{tab:promoter-elements} reports the resulting effect sizes.

The splice-shallow signature does not generalise to single-base
promoter elements. None of TSS, TATA, CAAT, or GC-box settles shallower
than the intronic baseline; TATA and GC-box settle measurably deeper
($d\!=\!+0.22$, $+0.29$ vs.\ intron). Under the GC-content-matched
control, only CAAT-box exhibits a small shallower-settling shift
($d\!=\!-0.12$), and GC-box remains robustly deep ($d\!=\!+0.31$). TATA
and TSS are statistically indistinguishable from their
composition-matched controls ($|d|\!\le\!0.05$), so the apparent depth
shift versus intron in those rows is largely a composition effect
rather than a residual-stream signature of motif detection. We read
this as direct support for arm~(b) of the bidirectional metric: at
promoter-element grammars, where Evo~2 must integrate CpG-island and
TF-binding-site context, $c(t)$ shifts deeper rather than shallower.
The splice-shallow signal of \S\ref{sec:shallowness} therefore reflects
a splice-specific motif-detection arm (arm~(a)) rather than a generic
short-motif phenomenon.

\begin{table}[!htbp]
  \caption{Canonical settling depth at chr22 promoter elements. ``vs.\ intron''
    is Cohen's $d$ against the chr22 intron baseline; ``vs.\ GC-matched
    control'' is Cohen's $d$ against a chr22 non-promoter control
    sampled in the same $\pm 5\%$ GC bin (1:1 matched). For reference,
    the chr22 splice-donor and splice-acceptor anchors yield
    $d\!=\!-0.367$ and $-0.351$ on the same intron baseline.}
  \label{tab:promoter-elements}
  \centering
  \small
  \begin{tabular}{@{}lcccc@{}}
    \toprule
    Promoter element & $n$ & $\overline{c}$ & $d$ vs.\ intron & $d$ vs.\ GC-matched ctrl \\
    \midrule
    TSS (Ensembl-canonical)    & $\phantom{00}{1{,}410}$ & $28.33$ & $+0.085$ & $+0.042$ \\
    TATA-box (MA0108.3, TBP)   & $\phantom{0}{13{,}560}$ & $29.14$ & $+0.219$ & $+0.027$ \\
    CAAT-box (MA0060, NFY)     & $\phantom{00}{4{,}392}$ & $28.11$ & $+0.049$ & $-0.121$ \\
    GC-box (MA0079, Sp1)       & $\phantom{00}{3{,}901}$ & $29.56$ & $+0.287$ & $+0.313$ \\
    \bottomrule
  \end{tabular}
\end{table}

\subsection{Per-Context Aggregation Consistency Check}
\label{app:agg-sanity}

Context-level summaries in the body use per-position $c(t)$ pooled over
windows: $\overline{c}_{\mathrm{ctx}}\!=\!|S_{\mathrm{ctx}}|^{-1}\sum_{t\in
S_{\mathrm{ctx}}} c(t)$. As a consistency check, we recomputed the
context-level ranking using per-region mean $c$ (one value per
annotated region rather than per position). The Spearman correlation
between the position-pooled and region-pooled per-context $\overline{c}$
vectors is $\ge 0.97$ on both chr22 and chr17, and the rank ordering
of the seven contexts is preserved. We report $d$ throughout because
it normalises by pooled within-context standard deviation and so
absorbs within-region variability by construction; both aggregation
schemes lead to the same qualitative ordering.

%%%%%%%%%%%%%%%%%%%%%%%%%%%%%%%%
% APPENDIX H — Extended Limitations (direction-only, composition,
% off-manifold shuffle)
%%%%%%%%%%%%%%%%%%%%%%%%%%%%%%%%
\section{Extended Limitations}
\label{app:extended-limitations}

This appendix expands the four-kind limitation summary in
\S\ref{sec:discussion}.

\subsection{Direction-only readout}
\label{app:lim-magnitude}
The cosine lens is scale-invariant, so anything carried in the residual
norm $\lVert h_\ell(t)\rVert$ is invisible to $c(t)$. Under a
superposition view~\cite{elhage2022toy,bricken2023monosemanticity}, the
norm reflects how strongly features are active, so two tokens that share
a direction but differ in norm get the same settling depth. We use
cosine on purpose: Evo~2's residual norms grow with depth, and a
raw-distance lens would confuse directional commitment with that
growth. The choice isolates the directional signal and gives up the
magnitude. Whether the discarded magnitude matters for settling is
testable: at layer~29, the spread of $\lVert h_{29}(t)\rVert$ across
tokens, and whether it adds anything to $c(t)$ for variant scoring
(App.~\ref{app:auroc-detail}), would measure what cosine leaves out.
Cosine is also not strictly magnitude-neutral, because a few outlier
dimensions can dominate it~\cite{timkey2021rogue}; confirming that
Evo~2's stream has no such dimensions, or standardising before the
lens, would put $c(t)$ on firmer ground. This is the same axis as the
absolute-versus-differential split in \S\ref{sec:variant-disruption}:
the absolute level saturates to a floor, which is fine for the
differential variant signal but lossy for the level-based settling
readout.

\subsection{Composition and the off-manifold shuffle}
\label{app:lim-composition}
The main standing confound is nucleotide composition. Local base
composition (GC content, dinucleotide frequencies) varies systematically
with annotation class, so a depth signal that aligns with composition
gradients risks being a composition artefact rather than a
representational signature. Our entropy decoupling
(\S\ref{sec:shallowness}, App.~\ref{app:entropy-decoup} and
\ref{app:restriction-of-range}) and GC-matched promoter controls
(App.~\ref{app:promoter-elements}) speak to this. The splice-donor
effect survives entropy regression, where $d$ strengthens from
$-0.452$ to $-0.583$ on the residual, and it is splice-specific rather
than a generic short-motif effect. The control is incomplete in two
places. The \S\ref{sec:motif-controls} flank shuffle keeps composition
fixed, so a composition-driven model would show a muted shuffle effect
rather than none. And App.~\ref{app:promoter-elements} shows the
confound is real for promoter elements, where TATA's $d=+0.22$ against
intron drops to $+0.03$ once GC is matched. We read the splice depth
signal as composition-robust and treat the promoter and flank-shuffle
contrasts as composition-entangled.

The flank-shuffle of \S\ref{sec:motif-controls} also has two readings
that the binary contrast cannot separate. Under \emph{reduced
integration}, depth tracks how much flanking grammar the model
integrates, so removing grammar settles $c$ earlier. Under \emph{early
simplification}, the model treats the shuffled flank as
composition-matched noise and commits to the lone GT. Both predict the
same direction; they differ only if early simplification is a
discontinuous response to noise rather than the smooth limit of graded
integration, and a graded-degradation series settles it: shuffle a
growing fraction of the flank, or shuffle at growing distance from the
junction. A monotone shift in $c$ supports reduced integration; a step
change supports a separate response to noise. The shuffle is also
off-manifold, since dinucleotide-shuffled flanks are out of
distribution for a model trained on natural genomes, so it reads
out-of-distribution behaviour alongside on-sequence integration. The
within-splice ordering in App.~\ref{app:splice-canonical} is the
on-manifold version of the same test, where non-canonical donors
settle earlier through a naturally tighter flank, so we treat it as
the primary evidence for context integration and the shuffle as its
controlled but out-of-distribution complement.

%%%%%%%%%%%%%%%%%%%%%%%%%%%%%%%%
% APPENDIX I — Reproducibility (last)
%%%%%%%%%%%%%%%%%%%%%%%%%%%%%%%%
\section{Data and Code Availability Statement}
\label{app:repro}

\paragraph{Code and data availability.}
The integrated analysis pipeline, source code, configuration files, and figure-generation scripts that underlie every numerical result in this paper are publicly available at \url{https://github.com/darejinn/gDTR}. All experiments use random seed $42$ for cross-validation splits and bootstrap resamples. The Evo~2 model is locked to \texttt{arcinstitute/evo2\_7b} at HF revision SHA \texttt{bda0089f92582d5baabf0f22d9fc85f3588f6b58} (weights MD5 \texttt{359ef88ccac2a62644035578de8a7db4}). Versioned external data sources comprise: GRCh38 primary assembly (UCSC, MD5 locked); GENCODE v44 GTF (per-chromosome filtered and persisted as a \texttt{gffutils} SQLite database); PhyloP 100-way (UCSC); ENCODE SCREEN v3 cCRE catalog (ELS subset); RepeatMasker hg38; GTEx v8 cis-eQTL pairs (\texttt{Whole\_Blood}, \texttt{Brain\_Cortex}, \texttt{Liver}, and \texttt{Lung} tissues, unioned); GWAS Catalog v1.0; and ClinVar archived release dated April 18, 2026. The reference software stack consists of \texttt{torch 2.4.1+cu124}, \texttt{evo2 0.3.0}, \texttt{vortex 1.0.8}, \texttt{transformer-engine 2.14.0}, \texttt{transformers 4.49.0}, \texttt{scipy 1.13}, and \texttt{scikit-learn 1.4}. Hardware: a single NVIDIA H200 (141\,GB). Total compute: approximately 20 GPU-hours end-to-end.

\paragraph{Data responsibility and ethics statement.}
The genomic coordinates and clinical variant datasets analyzed in this study were obtained from publicly available repositories under each repository's stated license. No restricted human-subject data, identifiable private health information, or individual-level genotype records were acquired or processed. All ClinVar records used here refer to aggregate submitter-classified variant calls without patient identifiers. The integrated pipeline complies with the data-governance frameworks of each upstream repository (UCSC, ENCODE, GTEx, GWAS Catalog, NCBI ClinVar) and uses an archival snapshot (ClinVar release: April 18, 2026) so that all numerical results in this paper remain exactly reproducible from that snapshot regardless of future repository updates.

\end{document}
